\chapter{Red Eye}

\section*{Definition}
Redness of the eye due to dilation of superficial blood vessels of the conjunctiva, episclera, or sclera, ranging from benign conjunctival injection to sight-threatening ocular emergencies.

\section*{What Happens in Your Body}
Red eye results from widening of blood vessels of conjunctival or episcleral vessels triggered by inflammation, infection, allergy, trauma, or increased intraocular pressure. Pattern of injection helps localize pathology: diffuse conjunctival injection suggests conjunctivitis; circumcorneal (ciliary flush) suggests keratitis, iritis, or acute glaucoma; sectoral injection suggests episcleritis or scleritis.

\section*{Causes}
\begin{itemize}
  \item Conjunctivitis (viral, bacterial, allergic) — most common
  \item Subconjunctival hemorrhage (benign, spontaneous)
  \item Corneal abrasion or foreign body
  \item Keratitis (infectious, contact lens-related)
  \item Acute angle-closure glaucoma (emergency)
  \item Uveitis or iritis
  \item Episcleritis or scleritis
  \item Dry eye disease (keratoconjunctivitis sicca)
  \item Trauma (corneal laceration, globe rupture)
  \item Orbital cellulitis (with proptosis, fever, pain with eye movement)
\end{itemize}

\section*{When to See a Doctor}
See your doctor if:
\begin{itemize}
  \item Symptoms are severe or getting worse over time
  \item Red eye with severe pain, vision loss, photophobia, halos around lights, corneal opacity, or trauma
  \item You have other concerning symptoms such as fever, weight loss, or bleeding
\end{itemize}

\section*{Related Symptoms}
\begin{itemize}
  \item Eye pain
  \item Vision blurred
  \item Photophobia
  \item Eye discharge
  \item Headache
\end{itemize}

\textbf{Important:} If this symptom persists, your doctor will check for serious underlying causes.

\section*{Self-Care / Home Management}
\begin{itemize}
  \item Rest and avoid activities that make it worse.
  \item Maintain adequate hydration and a balanced diet.
  \item Over-the-counter remedies may provide symptomatic relief (consult a pharmacist or doctor).
  \item Monitor symptoms and seek professional advice if they persist or worsen.
  \item Avoid self-medicating with prescription medications without medical guidance.
\end{itemize}

\section*{How Your Doctor Will Evaluate You}
\begin{itemize}
  \item A thorough history and physical examination are the first steps in evaluation.
  \item Laboratory tests (blood work, urinalysis) may be ordered based on clinical suspicion.
  \item Imaging studies (X-ray, ultrasound, CT, MRI) may be indicated when structural pathology is suspected.
  \item Specialist referral is arranged when the diagnosis remains uncertain or requires specific expertise.
\end{itemize}

\section*{Treatment Options}
\begin{itemize}
  \item Treatment targets the underlying cause identified through diagnostic evaluation.
  \item Supportive care and symptom management are provided while awaiting definitive therapy.
  \item Medications, lifestyle modifications, and physical therapies may be prescribed as appropriate.
  \item Follow-up ensures treatment effectiveness and allows timely adjustments to the care plan.
\end{itemize}

\clearpage